\chapter{Lithium-Ion Batteries: Vertical Depth}
\label{ch:batteries}

If solar demonstrated the cost-curve war, lithium-ion batteries demonstrated something further: the capture of an entire vertical---from mine to refinery to cathode to cell to pack---such that even a competitor who builds a gigafactory outside China still builds it, in effect, \emph{inside} the Chinese supply chain.

\section{From Japan's invention to China's entry (1991--2011)}

Sony commercialized the lithium-ion battery in 1991, building on Asahi Kasei's and John Goodenough's chemistry; Japan then Korea (LG Chem, Samsung SDI) led for two decades~\cite{sony-1991}. China's entry came through consumer electronics: BYD, founded in Shenzhen in 1995 with RMB 2.5 million and twenty employees, was by 2002 the world's largest nickel-cadmium battery maker with 65\% of global production, supplying Motorola and Nokia~\cite{byd-history}. CATL's lineage is similar: Robin Zeng's ATL (1999) built lithium-polymer cells for phones (and was sold to Japan's TDK); in 2011 Zeng spun the EV battery business out as CATL in Ningde~\cite{catl-history}. The pattern---enter at the low-margin consumer end, accumulate manufacturing capability, then pivot into the strategic segment when the state creates demand---recurs in Part~II.

\section{The white list: protection by regulation, not tariff (2015--2019)}

The decisive policy instrument was not a tariff but a certification. From 2015, NEV purchase subsidies applied only to vehicles using batteries from an MIIT-approved ``white list'' of 57 manufacturers---all domestic; LG and Samsung, which had just built Chinese plants, were excluded~\cite{white-list}. For four years, the world's largest and fastest-growing EV market was reserved, in effect, for CATL, BYD, and their domestic rivals. The list was scrapped in June 2019---precisely when the champions no longer needed it. Subsidy design was also used to steer chemistry: density-linked subsidies (2017--2019) pushed nickel-based NCM when that was deemed strategic, then cost-linked criteria (2020) resurrected LFP~\cite{lfp-story}. Regulation as industrial policy, adjusted dynamically, with foreign firms as designated losers.

\section{The end state: shares and costs}

\begin{table}[htbp]
  \centering
  \caption{Global EV battery installations, full-year 2025 (SNE Research)~\cite{sne-2025}.}
  \label{tab:batt-shares}
  \begin{tabular}{lS[table-format=4.1]S[table-format=2.1]}
    \toprule
    Maker & {GWh} & {Share (\%)} \\
    \midrule
    CATL (CN) & 464.7 & 39.2 \\
    BYD (CN) & 194.8 & 16.4 \\
    LG Energy Solution (KR) & 108.8 & 9.2 \\
    SK On (KR) & {---} & 3.7 \\
    Panasonic (JP) & {---} & 3.7 \\
    Samsung SDI (KR) & {---} & 2.4 \\
    \midrule
    World total & 1187.0 & 100.0 \\
    All Chinese makers & {$\sim$820} & {$\sim$69} \\
    \bottomrule
  \end{tabular}
\end{table}

The 2025 scoreboard (Table~\ref{tab:batt-shares}): a 1.19\,TWh market, with CATL and BYD alone above 55\% and Chinese makers collectively at ~69\% of installations and $\sim$75\% of deployment; China holds over 80\% of the world's 4\,TWh of cell manufacturing capacity, against 6--7\% each for the EU and US~\cite{sne-2025,iea-gevo-2026}. Upstream the shares are deeper still: roughly 85\% of anode capacity, 82\% of electrolytes, 74\% of separators, 70\% of cathodes; over 90\% of graphite processing; two-thirds of lithium and cobalt refining; and, through Tsingshan and partners, control of three-quarters of Indonesia's nickel refining~\cite{eia-components,benchmark-graphite,nickel}.

\begin{figure}[htbp]
\centering
\begin{tikzpicture}
\begin{axis}[bookaxis, height=5.0cm,
  ymode=log,
  xlabel={Year}, ylabel={Pack price (\$/kWh, log)},
  title={Lithium-ion pack prices: a 93\% real decline in fifteen years},
  xmin=2009.5, xmax=2026, ymin=70, ymax=2000,
  ytick={100,200,500,1000,2000}, yticklabels={100,200,500,1000,2000},
  xtick={2010,2013,2016,2019,2022,2025},
  legend pos=north east,
]
\addplot[color=sOne, mark=*, mark size=1.5pt] coordinates {
  (2010,1474) (2011,1100) (2012,900) (2013,732) (2014,668) (2015,384)
  (2016,295) (2017,214) (2018,180) (2019,156) (2020,137) (2021,132)
  (2022,151) (2023,139) (2024,115) (2025,108)};
\addlegendentry{global average pack}
\node[font=\scriptsize\color{inkSec}, anchor=south west] at (axis cs:2010.2,1500) {\$1,474};
\node[font=\scriptsize\color{inkSec}, anchor=north east] at (axis cs:2025.6,98) {\$108};
\end{axis}
\end{tikzpicture}

\vspace{1mm}

\begin{tikzpicture}
\begin{axis}[bookbar, height=4.0cm, width=0.62\linewidth,
  ylabel={2025 pack price (\$/kWh)},
  title={The cost gap, 2025},
  symbolic x coords={China,N.\ America,Europe},
  xtick=data, ymin=0, ymax=155,
  nodes near coords, nodes near coords style={font=\scriptsize\color{inkSec}},
]
\addplot[fill=sOne] coordinates {(China,84) (N.\ America,121) (Europe,131)};
\end{axis}
\end{tikzpicture}
\caption[Battery pack prices and the regional cost gap]{The second cost-curve war. \emph{Above}: BloombergNEF's
volume-weighted average pack price, real terms, 2010--2025~\cite{bnef-2025}; the 2022 uptick is the lithium
spike. \emph{Below}: the 2025 regional gap---North American and European packs run 44\% and 56\% above
Chinese~\cite{bnef-2025}. All values [V].}
\label{fig:battery}
\end{figure}


The cost curve replicated solar's. BNEF's pack price series fell from \$1,474/kWh (2010) to \$108/kWh (2025), a 93\% real decline; the China-specific price is \$84/kWh against \$121 in North America and \$131 in Europe---a 44--56\% Western cost penalty at the pack level~\cite{bnef-2025}.

\section{The LFP counterattack: innovation where the West wasn't looking}

The most strategically instructive episode is lithium iron phosphate. LFP was invented in the West (Goodenough's patent, 1997) and dismissed by Western industry as an energy-density dead end; forecasts gave it 15--40\% of the 2030 market. Chinese engineering---BYD's cell-to-pack Blade architecture, CATL's equivalent---recovered the density penalty at pack level, and LFP's cost, safety, and cycle-life advantages did the rest: LFP passed 50\% of the \emph{global} EV market in 2024 and 55\% in 2025~\cite{lfp-story,iea-gevo-2026}. The West now licenses the technology it invented back from China: Ford's Michigan plant operates under CATL license, and a CATL--Stellantis JV is building a €4.1 billion LFP plant in Spain~\cite{ford-catl}. The lesson for AI: the ``inferior'' technology branch, industrialized ruthlessly at scale, beats the premium branch on system economics---a preview of the LPDDR-versus-HBM argument of Chapter~\ref{ch:compute}.

China is now ahead on the next branches too: CATL's Naxtra sodium-ion line entered mass production in December 2025 (175\,Wh/kg, 10,000 cycles, $-40^{\circ}$C operation); second-generation Shenxing LFP charges 520\,km in five minutes; BYD's platform charges at one megawatt~\cite{naxtra,shenxing,byd-1mw}. Solid state remains genuinely open, but the IEA's sober assessment is prototypes industry-wide, with BYD targeting 2027~\cite{iea-gevo-2026}.

\section{From protection to gatekeeping}

The arc completed in 2025: the infant industry became the gatekeeper. In July 2025 MOFCOM placed LFP/LMFP cathode process technology and lithium-extraction technology under export license; in October the controls extended to high-energy cells ($\geq$300\,Wh/kg), anode materials, and the manufacturing equipment itself~\cite{mofcom-battery-controls}. China now restricts the export of battery \emph{technology} exactly as the US restricts semiconductors---the mirror image the West did not anticipate when it assumed technology flowed only one way.

\section{The Western casualty list}

The counterfactual competitors did not fail for lack of capital. Northvolt---Europe's champion, \$15 billion raised from VW, BMW, Goldman---collapsed into the largest industrial bankruptcy in modern Swedish history (November 2024 Chapter~11; Swedish bankruptcy March 2025) having produced 79.8\,MWh in three quarters of 2023, roughly \emph{one five-thousandth} of CATL's annual output~\cite{northvolt}. A123 had failed the same way a decade earlier, its assets bought by Wanxiang~\cite{a123}. By mid-2025, 46\,GWh of announced US cell capacity had been cancelled; LGES posted its first quarterly loss in three years and cut capex 30\%; Panasonic cut its North American output projections~\cite{us-cancellations,lg-loss}. Grid storage, the next battlefield, is already decided: China deployed $\sim$65\% of global BESS GWh in 2025, and storage packs from Chinese LFP oversupply price at \$70/kWh~\cite{bess-2025,bnef-2025}.

\section{Lessons carried forward}

Batteries add three lessons to solar's four. \emph{Fifth}: regulatory market reservation (the white list) is more effective than tariffs and invisible to trade law. \emph{Sixth}: vertical integration to the mine means a foreign competitor's factory is a hostage, not a rival---its cathodes, graphite, and equipment are Chinese regardless of its flag. \emph{Seventh}: the dismissed-technology branch (LFP) is where cost-curve wars are won; watch what China industrializes, not what the West benchmarks.
